\chapter{Estado del arte}
\label{cap:estado}

\section{Corrección automática de respuestas abiertas}

La corrección automática de texto libre, conocida en la literatura como \emph{Automatic Short Answer Grading} (ASAG), es un problema abordado desde distintos enfoques \parencite{burrows2015}. Los sistemas más simples se basan en reglas, listas de palabras clave o expresiones regulares. Su principal ventaja es la facilidad de implementación y la interpretabilidad, ya que permiten justificar la nota en función de elementos observables. Sin embargo, presentan una limitación evidente: penalizan en exceso la variabilidad lingüística. Cuando un estudiante expresa una idea correcta utilizando sinónimos, reformulaciones o una estructura distinta de la esperada, estos sistemas pueden infravalorar la respuesta.

En el extremo opuesto se sitúan métodos basados en modelos estadísticos o en representaciones semánticas densas. La aparición de los \emph{word embeddings} \parencite{mikolov2013} y, más tarde, de los modelos basados en \emph{transformers} permitió capturar significado más allá de la coincidencia literal, situando cerca, en un espacio vectorial, palabras o frases de sentido parecido. Estos enfoques son especialmente útiles cuando la diversidad de redacción es elevada. El inconveniente es que suelen requerir más datos, más capacidad de ajuste y, en ocasiones, sacrifican interpretabilidad. En entornos académicos, la posibilidad de explicar por qué una respuesta ha recibido una determinada calificación sigue siendo un requisito importante, y no siempre es compatible con un modelo cuya decisión es difícil de auditar.

Más recientemente, el uso de grandes modelos de lenguaje como evaluadores —el paradigma del \emph{LLM-as-a-judge} \parencite{zheng2023}— ha mostrado resultados notables, ya que estos modelos pueden valorar una respuesta a partir de una rúbrica expresada en lenguaje natural. No obstante, plantean dos problemas para un uso evaluativo serio: su salida no es plenamente reproducible y su razonamiento no es del todo transparente. Por ello, en este trabajo el modelo de lenguaje no se utiliza como juez principal, sino como una segunda opinión acotada y opcional, manteniendo el núcleo determinista.

En la práctica, este tipo de soluciones no siempre funcionan tan bien como se describe en teoría, especialmente cuando las respuestas son muy breves o poco estructuradas. Por ello, muchos sistemas útiles en escenarios docentes combinan distintas señales en lugar de depender de una única métrica.

\section{Conjuntos de datos de referencia}

Un aspecto recurrente en la investigación sobre ASAG es la dificultad de disponer de respuestas reales de alumnos corregidas por profesores. Existen, sin embargo, varios conjuntos de datos que sí incluyen esa información y que se han convertido en referencia para comparar sistemas.

El conjunto de \textcite{mohler2011} recoge respuestas de estudiantes de un curso universitario de informática, con alrededor de 2.400 respuestas a 87 preguntas, cada una calificada por dos correctores humanos en una escala de 0 a 5. Es especialmente valioso porque proporciona, junto a cada pregunta, una respuesta de referencia del profesor. El conjunto ASAP-SAS, difundido a través de una competición pública, ofrece un volumen mayor (en torno a 17.000 respuestas) sobre temas de ciencias a nivel escolar. Por su parte, la tarea 7 de SemEval-2013 \parencite{dzikovska2013} reúne casi 10.000 respuestas en quince dominios científicos, etiquetadas no solo con una nota sino con categorías de corrección.

Estos conjuntos comparten una limitación para el caso español: están redactados en inglés y, en su mayoría, pertenecen a dominios técnicos. Los exámenes oficiales en español que sí son públicos —por ejemplo, las pruebas de acceso a la universidad— suelen distribuir el enunciado, los criterios de corrección y una solución modelo, pero rara vez incluyen respuestas reales de alumnos con su nota, que es justamente lo que se necesita para medir la concordancia con un corrector. Esta asimetría se ha tenido en cuenta al diseñar la validación de este trabajo, que combina un banco propio en español con una prueba sobre datos reales en inglés.

\section{OCR en documentos académicos}

El reconocimiento óptico de caracteres constituye una etapa clave cuando el texto que se desea evaluar no se encuentra en formato digital nativo. Su función es transformar una imagen en una representación textual susceptible de ser procesada posteriormente. Aunque la idea general es sencilla, el rendimiento del OCR depende fuertemente de la calidad de la imagen. En documentos académicos escaneados pueden aparecer problemas como sombras, inclinación, ruido, contraste deficiente o variaciones en el tamaño de letra. Además, en el caso de respuestas escritas a mano o con formatos poco uniformes, la extracción se vuelve más inestable.

Tesseract \parencite{smith2007} es una de las herramientas más extendidas en este ámbito por su madurez, disponibilidad y capacidad para integrarse en \emph{pipelines} personalizados. En la literatura se asume con frecuencia que su rendimiento mejora cuando la imagen ha sido preprocesada de forma adecuada: convertir a escala de grises, binarizar, eliminar ruido o reforzar el contraste suelen presentarse como pasos que influyen en la calidad del texto extraído. No obstante, el beneficio real de cada una de estas operaciones depende del tipo de imagen y no puede darse por supuesto, por lo que conviene validarlo de forma empírica sobre los datos concretos del problema. Como se verá, en este proyecto esa validación llevó a una conclusión contraria a la intuición habitual.

Por tanto, en un sistema como el planteado, el OCR no puede considerarse un módulo aislado. Su salida afecta a todas las fases siguientes: un error de lectura puede alterar una palabra clave, fragmentar una frase o incluso desestructurar la separación entre pregunta y respuesta. Esto obliga a estudiar con cuidado qué tratamiento de imagen conviene aplicar antes de la extracción, valorando su efecto real sobre la lectura en lugar de asumir que cualquier preprocesamiento mejora el resultado.

\subsection{Texto manuscrito frente a texto impreso}

No es lo mismo reconocer texto impreso que manuscrito. El texto impreso presenta formas de letra regulares y un espaciado uniforme, lo que permite a un motor como Tesseract alcanzar precisiones altas con poco esfuerzo. La escritura manuscrita, en cambio, introduce una variabilidad enorme: cada persona traza las letras de forma distinta, las une de maneras diferentes y mantiene espaciados irregulares. A esto se añaden, en un examen real, las tachaduras, las correcciones, las flechas y las anotaciones al margen. Por eso el reconocimiento de manuscrito suele requerir modelos específicos, y su fiabilidad sigue siendo notablemente inferior a la del impreso.

Esta diferencia condiciona el alcance del proyecto. El sistema se ha diseñado y probado fundamentalmente sobre texto claro, asumiendo que la calidad de la entrada es razonable. Reconocer manuscrito complejo de forma fiable es un problema en sí mismo, ajeno al núcleo de este trabajo, que es la evaluación de la respuesta una vez extraída. Aun así, la robustez del módulo semántico frente a pequeños errores de lectura —mediante la detección parcial y la tolerancia a variaciones— amortigua parte de la imperfección del OCR, sin pretender resolverla por completo.

\section{Evaluación semántica y limitaciones de la comparación literal}

La comparación exacta de palabras es útil como señal inicial, pero insuficiente como criterio único de corrección. Existen varias razones para ello. En primer lugar, el lenguaje natural admite múltiples formulaciones equivalentes: una respuesta puede ser conceptualmente correcta aunque no reutilice los términos exactos del enunciado o de la solución modelo. En segundo lugar, la presencia de ciertas palabras no garantiza que el alumno haya entendido la idea; puede haber coincidencia léxica sin coherencia real, e incluso lo contrario, una negación que invierte el sentido. En tercer lugar, las respuestas abiertas suelen admitir distintos niveles de detalle, y la relevancia de cada fragmento no es homogénea.

La evaluación semántica trata precisamente de superar estas limitaciones. En términos generales, busca aproximar el significado global de un texto y valorar si contiene la información relevante desde el punto de vista conceptual. En aplicaciones reales esto suele resolverse combinando criterios locales y globales: por un lado, comprobar si aparecen elementos esenciales; por otro, medir si la respuesta, en conjunto, se parece a una referencia válida.

Un punto poco tratado en los sistemas más simples, pero crítico en evaluación, es la \emph{polaridad}. Un detector de palabras clave que solo comprueba presencia no distingue entre ``la mitocondria produce ATP'' y ``la mitocondria no produce ATP'': ambas contienen los mismos términos. La detección de la negación es un problema clásico del procesamiento del lenguaje natural, y su alcance —hasta dónde se extiende el efecto de un ``no'' dentro de una frase— condiciona la fiabilidad de cualquier corrector basado en conceptos. Este trabajo aborda explícitamente esa cuestión.

\subsection{Comparación entre enfoques de corrección}

Durante el diseño del sistema se consideraron tres estrategias principales de evaluación: comparación literal de palabras, similitud global del texto y enfoque híbrido. La diferencia entre ellas no es solo de complejidad, sino de comportamiento ante casos reales. En preguntas abiertas cortas, la elección del enfoque condiciona directamente qué errores se cometen y qué tipo de respuestas se valoran de forma injusta.

La comparación literal de palabras es el enfoque más simple. Consiste en verificar si la respuesta del alumno contiene exactamente los términos esperados. Su principal ventaja es la interpretabilidad inmediata: resulta sencillo justificar que una palabra clave aparece o no, y su implementación es poco costosa. También funciona razonablemente bien en preguntas muy cerradas. Sin embargo, su limitación es importante. Si un alumno responde correctamente usando sinónimos o reformulaciones, el sistema puede no reconocer la validez de la respuesta. Y, al revés, una frase incorrecta como ``la mitocondria es el núcleo de la célula'' contiene un término del dominio, pero eso no implica comprensión válida.

La similitud global del texto intenta corregir parte de ese problema al comparar la respuesta del alumno con una referencia completa. Su ventaja principal es que tolera mejor ciertos cambios de orden, estilo o formulación. No obstante, también presenta limitaciones claras. Puede penalizar respuestas correctas pero breves, porque la distancia con una referencia más desarrollada aumenta aunque el contenido esencial esté presente; puede verse afectada por errores de OCR que alteran varias palabras; y puede sobrevalorar respuestas superficialmente parecidas a la referencia aunque omitan un concepto esencial.

El enfoque híbrido utilizado en este proyecto combina lo más útil de ambos planteamientos y añade un control explícito sobre la longitud. La mayor parte de la nota depende de conceptos clave ponderados, mientras que la similitud global actúa solo como apoyo contextual. De esta forma el sistema mantiene interpretabilidad y, al mismo tiempo, gana tolerancia ante reformulaciones válidas. Si el alumno utiliza sinónimos o una redacción distinta, el bloque conceptual puede seguir detectando la evidencia relevante y la similitud aporta una corrección moderada. A la vez, si una respuesta reutiliza palabras del tema pero establece relaciones incorrectas, la mera coincidencia léxica no basta para obtener una puntuación alta.

En consecuencia, el enfoque híbrido resulta más adecuado para este proyecto porque se ajusta mejor a las condiciones reales del problema: respuestas abiertas, entrada afectada por OCR y necesidad de justificar la nota de forma comprensible. No es el enfoque más simple, pero sí el que ofrece un mejor equilibrio entre robustez, interpretabilidad y capacidad para distinguir entre respuestas completas, parciales e incorrectas.

\section{Herramientas y plataformas existentes}

En el ámbito comercial y académico existen plataformas que abordan partes del problema. Los sistemas de gestión del aprendizaje más extendidos incorporan corrección automática de cuestionarios, pero casi siempre limitada a preguntas de tipo test, emparejamiento o respuesta corta con solución exacta; las preguntas de desarrollo se siguen corrigiendo a mano. Existen también servicios especializados en evaluación de ensayos —el llamado \emph{automated essay scoring}— que puntúan redacciones extensas, habitualmente entrenados sobre grandes colecciones de textos calificados y orientados al inglés; su objetivo, sin embargo, es distinto del de este trabajo, ya que valoran sobre todo la calidad de la escritura y no la cobertura de unos conceptos concretos.

Más recientemente han aparecido asistentes basados en modelos de lenguaje que prometen corregir respuestas abiertas a partir de una indicación en lenguaje natural. Aunque resultan atractivos por su flexibilidad, comparten los inconvenientes ya señalados: su decisión es difícil de auditar y su salida no es del todo reproducible, lo que complica su uso en contextos donde una nota debe poder justificarse y defenderse. Frente a este panorama, la propuesta de este trabajo se diferencia en dos aspectos: mantiene un núcleo determinista e interpretable orientado a preguntas conceptuales, y trata el modelo de lenguaje como un complemento acotado en lugar de como el corrector principal.

\section{Métricas de evaluación de la concordancia}

Medir si un corrector automático ``acierta'' no es trivial, porque la propia noción de acierto admite matices. Comparar la nota del sistema con la del profesor mediante el error absoluto medio (MAE) informa de la distancia en la escala de calificación, pero penaliza por igual un sistema bien ordenado pero desplazado en escala y un sistema desordenado. Por eso suele acompañarse de medidas de correlación.

La correlación de Pearson cuantifica la relación lineal entre dos series de notas, mientras que la de \textcite{spearman1904} mide la relación de \emph{orden}: hasta qué punto el sistema clasifica las respuestas de mejor a peor igual que el humano, con independencia de la escala. En corrección, esta última es especialmente pertinente, porque a un profesor le importa sobre todo que el mejor examen reciba más nota que el peor, más que el decimal exacto. Una propiedad útil de la correlación de Spearman es que es invariante ante transformaciones monótonas de las notas, lo que la hace robusta frente a diferencias de escala y, como se verá, idónea para razonar sobre la calibración. En la literatura de evaluación automática también es común el coeficiente kappa ponderado cuadrático (QWK), habitual en competiciones del área.

\section{Calibración de puntuaciones}

Un sistema puede ordenar bien las respuestas y, sin embargo, puntuar en una escala distinta de la de un corrector concreto —por ejemplo, ser sistemáticamente más severo—. Ajustar la salida de un modelo a una escala de referencia es un problema clásico conocido como calibración. La forma más simple es una transformación lineal, definida por una pendiente y un sesgo. Cuando la relación no es lineal pero sí monótona, la regresión isotónica \parencite{barlow1972} ofrece una alternativa no paramétrica que ajusta una función creciente por tramos minimizando el error cuadrático. Su interés en este contexto es doble: al ser monótona no altera el orden de las respuestas —y por tanto no degrada la correlación de Spearman— y permite corregir desfases de escala arbitrarios con pocos datos. Estas ideas son la base de la capa de calibración descrita en el capítulo de metodología.

\section{Justificación del enfoque adoptado}

Tras revisar las alternativas, se ha optado por una solución híbrida y modular. Esta decisión responde a varios motivos. En primer lugar, permite construir un sistema realista sin depender de grandes volúmenes de datos etiquetados, algo poco compatible con el alcance de un Trabajo de Fin de Grado. En segundo lugar, facilita la explicación del resultado ante un tribunal o un docente, ya que cada componente de la nota puede justificarse. En tercer lugar, deja abierta la puerta a mejoras futuras: el sistema puede sustituir o enriquecer módulos concretos sin rediseñar la arquitectura completa.

Desde el punto de vista académico, esta elección resulta apropiada porque demuestra comprensión del problema y no se limita a ensamblar una herramienta existente. El valor del proyecto no reside solo en hacer funcionar un OCR, sino en definir cómo evaluar de forma razonable una respuesta abierta cuando la entrada procede de una imagen y cuando la corrección debe ser automática, interpretable y técnicamente defendible. La elección de mantener un núcleo determinista, y reservar el modelo de lenguaje para una función de verificación claramente delimitada, es coherente con ese objetivo.
